\documentclass[11pt,letterpaper]{article}
\usepackage[margin=0.7in]{geometry}
\usepackage[T1]{fontenc}
\usepackage{lmodern}
\usepackage{amsmath,amssymb,booktabs,array,xcolor,listings,hyperref,fancyhdr}
\definecolor{accent}{RGB}{25,65,100}
\hypersetup{colorlinks=true,linkcolor=accent,urlcolor=accent}
\setlength{\parindent}{0pt}
\setlength{\parskip}{4pt}
\setlength{\headheight}{14pt}
\pagestyle{fancy}\fancyhf{}\fancyhead[L]{\small CSI2110 -- Study Notes}\fancyhead[R]{\small Fall 2026}\fancyfoot[C]{\thepage}
\lstset{basicstyle=\small\ttfamily,breaklines=true,columns=fullflexible,keepspaces=true,frame=single,rulecolor=\color{accent!25},backgroundcolor=\color{accent!3},showstringspaces=false,tabsize=2}
\newcommand{\sourcepages}[1]{\textit{Source: PDF pages #1.}}
\newcommand{\keyidea}[1]{\par\textbf{Key idea.} #1\par}
\newcommand{\cost}[1]{\ensuremath{O(#1)}}
\newcommand{\method}[1]{\texttt{#1}}

\title{\textbf{Lecture 1: Analysis of Algorithms}\\[3pt]\large CSI2110 -- Data Structures and Algorithms}
\author{Lecture study notes}\date{}\begin{document}\maketitle
\textbf{Source document:} \path{../lec01-analysis of algorithms.pdf}.
These notes synthesize the source and add worked explanations. Page references use PDF page numbers, which can differ from printed slide numbers. Pseudocode is normalized for readability and boundary cases.
\tableofcontents\newpage

\section{Algorithms and efficiency}
\sourcepages{4--14}
An algorithm is a finite, well-defined procedure for solving a problem. Evaluate \textbf{correctness}, \textbf{running time}, and \textbf{memory usage}. The input size $n$ measures the amount of data; for an array it is usually the number of entries.

Different inputs of the same size can have different costs. Linear search may find a target immediately, at the last index, or not at all. Its best-case time is $\Theta(1)$ and worst-case time is $\Theta(n)$.

Experimental analysis implements and times a program on varied inputs. It reveals practical performance, but depends on hardware, language, compiler, and the chosen inputs. Theoretical analysis counts operations independently of a particular machine.

\section{Pseudocode and the RAM model}
\sourcepages{15--19}
Pseudocode expresses an algorithm without unnecessary implementation details. Indentation defines blocks; assignment is written $\leftarrow$, while $=$ means equality.

The random access machine (RAM) model assumes constant-time memory access and constant-time primitive operations such as assignment, arithmetic on machine-sized values, comparisons, and return. A loop or a method call is not automatically one primitive operation: count the work it performs.

\subsection{Worked example: finding the maximum}
Assume $n\geq1$.
\par\noindent\begin{minipage}{\linewidth}
\begin{lstlisting}
arrayMax(data):
    currentMax = data[0]
    for i = 1 to data.length - 1:
        if data[i] > currentMax:
            currentMax = data[i]
    return currentMax
\end{lstlisting}
\end{minipage}\par
\sourcepages{19--31}
There are always $n-1$ element comparisons. The current maximum is assigned once initially and then at most $n-1$ more times. The lecture's particular counting convention gives $4n+3$ primitive operations in the best case and $6n+1$ in the worst case. Exact constants depend on the counting convention; both cases take $\Theta(n)$ time.

\textbf{Invariant:} before iteration $i$, \method{currentMax} is the maximum of entries $0,\ldots,i-1$. Comparing entry $i$ maintains the invariant. At termination all entries have been considered.

\section{Asymptotic notation}
\sourcepages{33--54,62--68}
Asymptotic analysis describes growth for sufficiently large $n$. Assume the compared functions are eventually nonnegative and $g(n)$ is eventually positive.
\begin{align*}
f\in O(g)&\iff \exists c>0, n_0:\quad f(n)\leq c g(n)\quad\forall n\geq n_0,\\
f\in\Omega(g)&\iff \exists c>0, n_0:\quad f(n)\geq c g(n)\quad\forall n\geq n_0,\\
f\in\Theta(g)&\iff \exists c_0,c_1>0, n_0:\quad
c_0g(n)\leq f(n)\leq c_1g(n)\quad\forall n\geq n_0.
\end{align*}
\begin{itemize}
\item $O$ is an upper bound, $\Omega$ a lower bound, and $\Theta$ a tight bound.
\item $c$ absorbs constant factors; $n_0$ excludes a finite initial range.
\item $c$ and $n_0$ must be constants independent of $n$. They need not be the smallest possible values.
\item $\Theta(g)=O(g)\cap\Omega(g)$ and $f\in O(g)$ iff $g\in\Omega(f)$.
\end{itemize}
\keyidea{$O$ does not mean ``worst case'' and $\Omega$ does not mean ``best case.'' First choose the running-time function (best, worst, or average case), then bound that function.}

\subsection{Proofs by definition}
For $n\geq1$,
\[
25n+2110\leq25n+2110n=2135n.
\]
Thus $25n+2110\in O(n)$ with $c=2135,n_0=1$.

For $n\geq1$,
\[
3n^3+20n^2+5\leq28n^3.
\]
Hence the polynomial is $O(n^3)$. It is also $\Omega(n^3)$ because it is at least $3n^3$, so it is $\Theta(n^3)$.

For $n\geq1$,
\[
20n^2\leq20n^2+4n\leq24n^2,
\]
proving $\Theta(n^2)$ with $c_0=20,c_1=24,n_0=1$.

For base-2 logarithms and $n\geq2$, $\log_2 n\geq1$, so
\[
n\log_2 n\leq n\log_2 n+204n\leq205n\log_2 n.
\]
This proves $\Theta(n\log n)$. The threshold $n_0=2$ is needed for this upper bound; $n=1$ does not satisfy it.

\subsection{Rules and growth order}
For nonnegative functions:
\begin{itemize}
\item Reflexivity: $f\in O(f)$; transitivity: $f\in O(g),g\in O(h)\Rightarrow f\in O(h)$.
\item If $f,g\in O(h)$, then $f+g\in O(h)$.
\item If $f_1\in O(f_2)$ and $g_1\in O(g_2)$, then $f_1g_1\in O(f_2g_2)$.
\item $\max\{f,g\}\leq f+g\leq2\max\{f,g\}$, so $f+g\in\Theta(\max\{f,g\})$.
\end{itemize}
Drop constant factors and lower-order terms when finding a tight class. An eventually nonnegative degree-$d$ polynomial with positive leading coefficient is $\Theta(n^d)$. Using absolute coefficient values gives a valid upper-bound proof even when some coefficients are negative.

Common growth rates, in increasing order, are
\[
1,\quad\log n,\quad n,\quad n\log n,\quad n^2,\quad n^3,\quad
2^n,\quad n2^n,\quad e^n,\quad n!.
\]
A fixed logarithm base greater than one changes only a constant factor. A fixed exponential base \emph{can} change the growth class: $3^n\notin O(2^n)$.

\subsection{Lecture classification exercise}
\sourcepages{66}
\begin{center}\begin{tabular}{@{}lccc@{}}\toprule
Function & $O(n^3)$ & $\Omega(n^3)$ & $\Theta(n^3)$\\\midrule
$77n^2$ & Yes & No & No\\
$6n^3+n^2\log_2 n$ & Yes & Yes & Yes\\
$11n^4$ & No & Yes & No\\
$n^2$ if even; $n^4$ if odd & No & No & No\\\bottomrule
\end{tabular}\end{center}
The last function violates the upper bound on arbitrarily large odd inputs and the lower bound on arbitrarily large even inputs. Asymptotic definitions require the inequality for \emph{all} sufficiently large inputs.

\section{Analyzing loops}
\sourcepages{56--60}
Count iterations, multiply by the work per iteration, and sum costs when bounds depend on an outer variable.
\begin{center}\begin{tabular}{@{}ll@{}}\toprule
Pattern, with constant work in the body & Tight time\\\midrule
$i\leftarrow i+1$, up to $n$ & $\Theta(n)$\\
$i\leftarrow i+2$, up to $n$ & $\Theta(n)$\\
$i\leftarrow2i$, starting at $1$, up to $n$ & $\Theta(\log n)$\\
A fixed 1001 iterations & $\Theta(1)$\\
Two nested loops, each $n$ iterations & $\Theta(n^2)$\\
Outer $i=0,\ldots,n-1$; inner $j=0,\ldots,i$ & $\Theta(n^2)$\\\bottomrule
\end{tabular}\end{center}
The triangular case has $\sum_{i=0}^{n-1}(i+1)=n(n+1)/2$ iterations. Nested loops do not always multiply to $n^2$; use their actual bounds. The fixed-loop array example requires a large enough array for its indices to be legal.

\section{Prefix averages: avoiding repeated work}
\sourcepages{70--76}
Given $X[0],\ldots,X[n-1]$, define
\[
A[i]=\frac{\sum_{j=0}^{i}X[j]}{i+1}.
\]
A direct algorithm recomputes each prefix sum:
\par\noindent\begin{minipage}{\linewidth}
\begin{lstlisting}
for i = 0 to n - 1:
    s = 0
    for j = 0 to i:
        s = s + X[j]
    A[i] = s / (i + 1)
\end{lstlisting}
\end{minipage}\par
It performs $1+2+\cdots+n=\Theta(n^2)$ additions. Maintain a running sum instead:
\par\noindent\begin{minipage}{\linewidth}
\begin{lstlisting}
s = 0
for i = 0 to n - 1:
    s = s + X[i]
    A[i] = s / (i + 1)
\end{lstlisting}
\end{minipage}\par
This takes $\Theta(n)$ time. Both allocate $\Theta(n)$ output space and use $O(1)$ additional working space. Use real-number division for averages.

For $X=[5,13,4,8]$, the running sums are $5,18,22,30$, giving averages $5,9,22/3,7.5$.

\section{Prefix averages at geometric steps}
\sourcepages{77--90}
Only return averages for prefix lengths $1,2,4,\ldots,2^k$, where $k=\lfloor\log_2 n\rfloor$ and $n\geq1$:
\[
A[i]=\frac{\sum_{j=0}^{2^i-1}X[j]}{2^i},\qquad 0\leq i\leq k.
\]
Recomputing these sums costs
\[
\sum_{i=0}^k2^i=2^{k+1}-1=\Theta(n),
\]
not $\Theta(n\log n)$. A single scan with a running sum and a doubling next-prefix-length variable also takes $\Theta(n)$ time. Both approaches have $\Theta(\log n)$ output space and $O(1)$ working space. Thus avoiding recomputation here improves constants but not the asymptotic class.
\par\noindent\begin{minipage}{\linewidth}
\begin{lstlisting}
s = 0; nextLength = 1; i = 0
for j = 0 to n - 1:
    s = s + X[j]
    if j + 1 == nextLength:
        A[i] = s / nextLength
        i = i + 1
        nextLength = 2 * nextLength
\end{lstlisting}
\end{minipage}\par

\section{Mathematical toolkit}
\sourcepages{51--53,79,92--94}
\begin{align*}
a^{b+c}&=a^ba^c,& (a^b)^c&=a^{bc},& a^b/a^c&=a^{b-c},\\
\log_b(xy)&=\log_bx+\log_by,& \log_b(x/y)&=\log_bx-\log_by,\\
\log_b(x^a)&=a\log_bx,& \log_bx&=\frac{\log_cx}{\log_cb}.
\end{align*}
Here logarithm arguments are positive and bases are positive and unequal to one. $\lfloor x\rfloor$ is the largest integer at most $x$; $\lceil x\rceil$ is the smallest integer at least $x$.
\[
\sum_{i=1}^n i=\frac{n(n+1)}2,\qquad
\sum_{i=0}^n id=d\frac{n(n+1)}2,\qquad
\sum_{i=0}^n r^i=\frac{r^{n+1}-1}{r-1}\quad(r\ne1).
\]
For fixed $r>1$, the geometric sum is $\Theta(r^n)$; for $r=1$ it is $n+1$; for fixed $0<r<1$ it is $\Theta(1)$. The arithmetic sum is $\Theta(n^2)$ when $d>0$ is fixed.

Optional: Stirling's approximation is $n!\sim\sqrt{2\pi n}(n/e)^n$.

\paragraph{Study checklist}
Be able to give explicit constants for an asymptotic proof, distinguish a loose upper bound from a tight bound, count triangular and geometric loops, and explain why running sums improve ordinary prefix averages but not geometric prefix averages.

\end{document}
